\section{El Ppio. de Acotación Uniforme}

\begin{ejercicio}
Probar que, en el espacio de Banach $C\left[0, 1\right]$, las funciones que son derivables en el punto $\nicefrac{1}{2}$ forman un conjunto de primera categoría.
\end{ejercicio}

\begin{ejercicio}
Sea $y \in \bb{K}^{\bb{N}}$ una sucesión de escalares y sea $1 < p < \infty$. Supongamos que, para todo $x \in \ell_p$, la serie
\[
\sum_{n \geq 1} x(n)y(n)
\]
está acotada. Probar que $y \in \ell_{p^*}$.
\end{ejercicio}

\begin{ejercicio}
Sean $X$ e $Y$ espacios de Banach y sea $T: X \longrightarrow Y$ un operador lineal. Probar que $T$ es un monomorfismo si, y sólo si, existen $\alpha, \beta \in \bb{R}^+$ tales que:
\[
\alpha \|x\| \leq \|T(x)\| \leq \beta \|x\| \quad \forall x \in X
\]
\end{ejercicio}

\begin{ejercicio}
Probar que el conjunto de los monomorfismos entre dos espacios de Banach es abierto en el espacio de todos los operadores lineales y continuos.
\end{ejercicio}

\begin{ejercicio}
Sean $X$ e $Y$ espacios de Banach y $T \in L(X, Y)$. Probar que las siguientes afirmaciones son equivalentes:
\begin{enumerate}[label=\textbf{\alph*)}]
    \item Existe $S \in L(Y, X)$ tal que $S(T(x)) = x$ para todo $x \in X$.
    \item $T$ es inyectivo y $T(X)$ es un subespacio complementado de $Y$.
\end{enumerate}
\end{ejercicio}

\begin{ejercicio}
Sean $X$ un espacio de Banach, $Y$ un espacio normado y $T: X \longrightarrow Y$ una aplicación lineal. Se define una nueva norma en $X$ mediante la expresión $\|x\|_1 = \|x\| + \|T(x)\|$ para todo $x \in X$. Probar que las siguientes afirmaciones son equivalentes:
\begin{enumerate}[label=\textbf{\alph*)}]
    \item $T$ es continua.
    \item $\|\cdot\|$ y $\|\cdot\|_1$ son normas equivalentes.
    \item $\|\cdot\|_1$ es una norma completa en $X$.
\end{enumerate}
\end{ejercicio}

\begin{ejercicio}
Probar que no existe ninguna sucesión $u = \{u_n\}$ en $\bb{K}$ tal que, para toda sucesión $x = \{x_n\}$ en $\bb{K}$, se verifica que:
\[
x \in \ell_1 \sii \{x_n u_n\} \text{ está acotada.}
\]
[Sugerencia: Supuesta la existencia de una tal sucesión $u = \{u_n\}$, nótese que se puede suponer que todos sus términos son no nulos, y considérese para cada sucesión $x = \{x_n\}$ la sucesión $x' = \left\{\frac{x_n}{u_n}\right\}$, y demuéstrese que la aplicación $x \mapsto x'$ sería un isomorfismo topológico de $\ell_\infty$ sobre $\ell_1$].
\end{ejercicio}

\begin{ejercicio}
Sea $X$ un espacio de Banach real. Dada una aplicación lineal $T: X \longrightarrow L_1\left(\left[0, 1\right]\right)$ se considera, para cada $A \subseteq \left[0, 1\right]$ medible, el funcional lineal $T_A: X \longrightarrow \bb{R}$ definido por
\[
T_A(x) = \int_A T(x)\,dx
\]
Probar que si $T_A \in X^*$ para todo $A \subseteq \left[0, 1\right]$ medible, entonces $T$ es continua.
\end{ejercicio}

\begin{ejercicio}
Sean $X$ un espacio de Banach real y $T: X \longrightarrow C\left(\left[0, 1\right], \bb{R}\right)$ una aplicación lineal. Se considera, para cada $n \in \bb{N} \cup \{0\}$, el funcional lineal $T_n: X \longrightarrow \bb{R}$ definido por
\[
T_n(x) = \int_{0}^1 t^n T(x)(t)\,dt
\]
Probar que si $T_n \in X^*$ para todo $n \in \bb{N} \cup \{0\}$, entonces $T$ es continua.
\end{ejercicio}

\begin{ejercicio}
Sean $(E, d)$ un espacio métrico y $X$ un espacio normado. Probar que una función $F: E \longrightarrow X$ es lipschitziana si, y sólo si, para todo $x^* \in X^*$, la función $x^* \circ F: E \longrightarrow \bb{K}$ lo es.
\end{ejercicio}

\begin{ejercicio}
Sea $\cc{H}$ el espacio vectorial de las funciones enteras de variable compleja.
\begin{enumerate}[label=\textbf{\alph*)}]
    \item Probar que la fórmula
    \[
    \|f\| = \max\{|f(z)| : |z| = 1\}
    \]
    define una norma en $\cc{H}$.
    \item Probar que la aplicación $F: \cc{H} \longrightarrow \cc{H}$ definida por
    \[
    \left[F(f)\right](z) = f\left(\frac{z}{2}\right) \quad (f \in \cc{H}, \; z \in \bb{C})
    \]
    es una biyección lineal y continua. (Indicación: Utilizar el principio del módulo máximo).
    \item Probar que $F^{-1}$ es discontinua.
    \item ¿Es $\cc{H}$ un espacio de Banach?
\end{enumerate}
\end{ejercicio}